
We trained ResNet-50 on Waterbirds and found that --- without any group labels --- the curvature of the loss surface around each training sample reveals minority group membership with AUROC 0.85--0.90 in 4/5 random seeds (mean AUROC 0.89 across the four passing seeds). When investigating the mechanism, however, the expected mechanistic driver was absent: by the time the curvature signal peaks, minority samples are classified with training confidence $\geq 0.97$, indistinguishable from majority samples. The signal is robust; the original mechanistic explanation is not.

This puzzle sits at the intersection of two long-standing problems in robust machine learning. The first is spurious correlations: models trained with ERM on real-world datasets learn to exploit incidental statistical associations between input features and labels. On Waterbirds \citep{sagawa2019distributionally}, a ResNet-50 trained with ERM achieves approximately 97\% average accuracy on training samples but only approximately 72\% worst-group accuracy on the test set --- minority groups (e.g., landbirds on water backgrounds) suffer because the model relies on background rather than bird shape. The second problem is annotation burden: the most effective mitigation methods, such as deep feature reweighting (DFR) \citep{kirichenko2022last}, require a balanced held-out validation set with explicit group labels --- a resource that is unavailable in many practical settings.

Annotation-free alternatives --- Just Train Twice (JTT) \citep{liu2021just}, SELF \citep{labonte2023towards} --- proxy minority membership through first-order signals: which samples are misclassified, or which have the highest loss. These methods improve over ERM without group labels, but they are insensitive to how the model encodes spurious features in its loss landscape geometry.

This work asks whether the second-order geometry of the loss surface carries distinct information about minority group membership. Specifically, we measure the per-sample Hessian trace of the last fully-connected layer --- computed via K=50 Hutchinson estimation using torch.func vmap+vjp --- at six training checkpoints ($t\in {0,1,5,10,20,50})$ over five random seeds, and ask whether this trajectory discriminates minority from majority samples.

The measurement protocol is motivated by two theoretical results. LaBonte et al. [2024] show that minority group covariance matrices have larger spectral norm than majority groups within the same class --- a group-level prediction of systematic feature-norm inequality ($\|x_{i}\|^{2})$ that would elevate per-sample Hessian traces for minority samples, given the decomposition Tr($H_i^{fc}$) ∝ $\|x_{i}\|^2$ p\_i($1-p_{i})$. LaBonte and Muthukumar [2026] prove that SGD on XOR spurious models learns the spurious feature first (Phase I) before the core feature (Phase II), predicting a transient differential in loss landscape curvature between majority (which gain confidence rapidly via the spurious shortcut) and minority (which do not).

The existence experiment (H-E3) confirms the signal: 4/5 seeds achieve $AUROC\geq 0.85$ for minority membership prediction at t* (the epoch maximizing the trace ratio R(t) = mean\_minority\_trace / mean\_majority\_trace), with epoch-0 AUROC<0.70 in all seeds. The Hutchinson estimator is stable (CV<3\% for all seeds), establishing K=50 as sufficient for the last-fc layer of ResNet-50.

The mechanism experiment (H-M1) produces the central puzzle. The original mechanistic explanation --- that minority samples remain near the decision boundary ($p_i\in [0.3,0.7]$) throughout Phase I training, keeping p\_i($1-p_{i})$ large and thereby elevating Tr($H_i^{fc}$) --- is directly falsified. At t\textit{, minority training-set confidence is 0.9678--0.9999 across all seeds: fully saturated, not boundary-straddling. A transient differential does exist at epoch 1 (seed 1: p\_min=0.827 vs p\_maj=0.972), consistent with Phase I theory, but it disappears well before t} in all seeds.

The trace signal at t\textit{ must therefore arise from the $\|x_{i}\|^2$ feature-norm channel --- systematic differences in the penultimate-layer representation norms between minority and majority samples, consistent with LaBonte et al. [2024]'s spectral imbalance finding. This interpretation is a candidate, not a verified claim; directly measuring penultimate-layer norms at t} is the primary direction for future mechanistic investigation.

\textbf{Contributions:}

\begin{enumerate}
\item \textbf{Empirical existence result:} First experimental evidence that per-sample last-fc Hessian trace (K=50 Hutchinson via vmap+vjp) achieves $AUROC\geq 0.85$ for minority membership detection in ERM-trained ResNet-50 on Waterbirds across 5 random seeds, with epoch-0 control confirming ERM emergence (Section 5.1).
\end{enumerate}

\begin{enumerate}
\item \textbf{Principled falsification of the confidence-differential mechanism:} Direct refutation of the sustained decision-boundary hypothesis ($p_{\text{minority}}\in [0.3,0.7]$ at t*) on training-set data, with identification of the feature-norm channel ($\|x_{i}\|^{2})$ as the likely surviving driver (Section 5.2, Section 6.2).
\end{enumerate}

\begin{enumerate}
\item \textbf{Practical confirmation:} K=50 Hutchinson is the minimum stable setting for last-fc ResNet-50 (CV<3\%), and the epoch-0 control serves as a reliable pretrained-artifact detector (Section 4.4).
\end{enumerate}
